\section{Introduction}
\label{sec:intro}

Generalist vision-language-action (VLA) policies~\cite{brohan2023rt2,kim2024openvla,black2024pi0,pi2025pi05} now solve
the standard LIBERO suites~\cite{liu2023libero} almost perfectly.
The same policies fail in a strikingly simple way when the scene is rearranged.
The \emph{swap} perturbation of LIBERO-PRO~\cite{zhou2025liberopro} keeps every task and only lets the target trade places
with another object; under it, \pizf drops from 99\% to 19\% on LIBERO-Object and from 97.5\% to 40\% on LIBERO-Spatial.
Watching the rollouts shows why: the arm drives to where the target used to be.
The policy behaves as if the instruction retrieved a memorized trajectory~\cite{chen2026ect}.

A natural reading is that the policy does not \emph{perceive} where things are, and a large body of work adds
perception: depth, point clouds, 3D features, or an explicit 3D goal point~\cite{zhen20243dvla,qu2025spatialvla,tsai2026point}.
We test this reading directly and find it wrong for \pizf.
Because our policies run in a simulator, we can stop any rollout at any state, change exactly one thing (move the target,
swap it with another object, move the container while the object is carried, rename the target in the instruction),
re-render, and ask the policy again.
Comparing the two predictions measures how much the action \emph{uses} the change; probing the hidden states before and
after the change measures how well the change is \emph{encoded}.
The two disagree (\cref{fig:teaser}).
The displacement is linearly decodable from the action tokens of \pizf ($R^2$ 0.55--0.86), but at the moments that decide
the task, when the hand is still on its way to the target or carries the object to the container, the predicted action
follows the displacement by only 2--25\%.
\pizf \emph{knows where} the target is and does not \emph{go there}.

This distinction predicts which remedies can work, and \cref{sec:theory} makes the prediction precise: on fixed layouts
the data cannot tell a policy that uses the target's position from one that replays the memorized one.
If the information is already present, adding more of it should not help, and it does not: re-implementing the injection
of a privileged 3D goal point into the action expert~\cite{tsai2026point} in our setting leaves swap success at the base
level, and the injected channel is learned as a constant offset that ignores the point.
Equivariant counterfactual training (\ect)~\cite{chen2026ect}, which pairs demonstrations with geometrically transformed
replays, diagnoses the same memorization but, in a matched setting, leaves the measured use where it was.
What does change use is supervision \emph{at the decision states themselves}: building, at a state of the policy's own
rollout, a world in which the memorized action is wrong (the target or the container has moved) and supervising the
correct action there.
We call this decision-time counterfactual supervision (\ours).
It raises use from 0.02--0.06 to 0.39--0.68 while the encoding does not change: it teaches the policy to act on what it
already represents rather than to see better.

The same instrument tells us \emph{where} the change happens.
Restricting the trainable parameters to different parts of \pizf, with the same counterfactual data, shows that the
action expert, its readout layer, or the keys and values the action expert reads from the VLM are not where the fix lives;
adapting the VLM backbone alone, with the action expert frozen, recovers about 90\% of the swap gain, and most of that
comes from its layers 9--17.

\paragraph{Contributions.}
\begin{itemize}
\item A decision-time intervention protocol that separates \emph{use} from \emph{encoding} in a VLA, per task phase, and
the finding that \pizf encodes object and container displacements it does not act on (\cref{sec:probing}).
\item An identifiability account (\cref{sec:theory}): when layouts are fixed, the training risk is indifferent to whether
a policy uses the positions it encodes, so adding information (\eg a 3D goal point) cannot help, while counterfactuals
that move a single decision variable at decision states, with labels that depend on it, identify use. Five predictions of
this account hold in our experiments, including that privileged 3D goal injection is ignored (\cref{sec:exp:mech}).
\item A localization of the change that does fix use: the VLM's mid-to-late layers, not the action pathway
(\cref{sec:exp:local}).
\item Decision-time counterfactual supervision with validity-gated privileged labels, which improves LIBERO-PRO swap by
31.3 points (LIBERO-Object) and 14.8 points (LIBERO-Spatial) over \ect under the same teacher and budget
(\cref{sec:exp:main}), together with an account of its costs and of an evaluation detail of LIBERO, the per-reset
re-placement of furniture, that changes per-task results by up to 25 points (\cref{sec:exp:limits}).
\end{itemize}
